\subsection*{C++ Interop: Device Functions, Inline Code, and the Preprocessor}

A kernel may need a reduction, a sorting routine, or an instruction that Croqtile does not express directly. Existing device C++ can supply that computation while Croqtile continues to arrange its inputs, outputs, and parallel scope. The interface needs to preserve the assumptions on both sides of the call. The reason is always the same: no matter how expressive the abstraction, some machine feature, some legacy library, some performance-critical intrinsic lives outside the language's domain.

Croqtile's interoperability story has three parts:

\begin{enumerate}[leftmargin=1.6em]
\item \textbf{\texttt{\_\_device\_\_} functions} --- device functions that coexist with \texttt{\_\_co\_\_} functions in the same \texttt{.co} file. The compiler passes them through unchanged.
\item \textbf{\texttt{\_\_cpp\_\_}} --- inject verbatim C++ into the generated code for target-specific instructions the DSL does not emit.
\item \textbf{The preprocessor} --- \texttt{\#define}, \texttt{\#if}, \texttt{\#ifdef}, \texttt{\#error} for compile-time configuration, the same role \texttt{\#define} plays in C/C++ codebases.
\end{enumerate}

The preceding sections introduced tiles, parallelism, MMA, pipelines, and events. That is the Croqtile you want to live in. This part of the guide is about the times you need to step outside.
